\documentclass[11pt,a4paper]{article}
\usepackage[utf8]{inputenc}
\usepackage[T1]{fontenc}
\usepackage[spanish]{babel}
\usepackage{lmodern}
\usepackage[margin=2.2cm]{geometry}
\usepackage{graphicx}
\usepackage{booktabs}
\usepackage{tabularx}
\usepackage{hyperref}
\usepackage[backend=biber,style=ieee]{biblatex}
\addbibresource{referencias.bib}

\begin{document}

\begin{center}
    \textbf{\Large UNIVERSIDAD TÉCNICA ESTATAL DE QUEVEDO}\\
    \textbf{Facultad de Ciencias de la Computación y Diseño Digital}\\
    \textbf{Carrera de Ingeniería en Software -- Paralelo A}\\[0.4cm]
    \textbf{\large Notas de Presentación: WebSockets, TCP y UDP aplicados a JobFinder}\\
    \small Actividad: GA-SUM-02 / EXP-1 (Exposición 1.er Corte) -- Equipo EJPZ
\end{center}

\vspace{0.2cm}
\noindent\textbf{Integrantes:} España Vera Kenny Isaías, Jaramillo Acosta Luis Jeremy, Paca Pilataza Luis Alberto, Zambrano Ronquillo Heiner Ariel.

\section{Guion Desarrollado por Turnos de Exposición (12 Minutos)}

\subsection{Bloque 1: Fundamentos de Transporte y Presentación del Sistema (Minutos 0:00--3:00) -- Expositor: Kenny Isaías España Vera}
En la capa de transporte, TCP establece una sesión mediante saludo de tres vías y asegura entrega ordenada mediante acuses de recibo (ACK) y retransmisión, mientras que UDP envía datagramas independientes sin conexión ni garantía de recepción. El protocolo WebSocket (RFC 6455) opera sobre TCP iniciando con una negociación HTTP (\texttt{Upgrade: websocket}) que responde \texttt{101 Switching Protocols}, habilitando un canal \textit{full-duplex} persistente sin la sobrecarga de cabeceras del \textit{polling} tradicional.

\textbf{Presentación de JobFinder (Punto 1):} JobFinder resuelve la baja inserción laboral institucional de los graduados de la UTEQ, donde el seguimiento institucional reporta que solo el 51,50\,\% ejerce y apenas el 8,41\,\% consigue empleo mediante bolsas universitarias. El sistema atiende a graduados, empresas colaboradoras y administradores. \textit{Práctica 1:} ejecución de datagramas UDP en consola frente a la inspección del \textit{handshake} HTTP 101 en Chrome DevTools.

\subsection{Bloque 2: Arquitectura y Modelos de Comunicación (Minutos 3:00--6:00) -- Expositor: Luis Jeremy Jaramillo Acosta}
\textbf{Arquitectura (Punto 2):} Elaborada en \texttt{draw.io} en tres niveles sobre TCP: (1) clientes web en Angular (\texttt{:4200}); (2) backend en Spring Boot (\texttt{:8080}) compuesto por API REST (\texttt{/api/**}), módulo WebSocket y notificaciones (actualmente STOMP/SockJS en \texttt{/ws-notificaciones} y extendido en nuestra práctica con \texttt{ws://} puro \textit{full-duplex}) e integración con la API externa JSearch; y (3) persistencia en PostgreSQL (\texttt{:5432}, base de datos \texttt{BolsaEmpleoNew} con auditoría en \texttt{seguridad.auditoria}) junto a correo SMTP (\texttt{smtp.gmail.com:587}).

\textbf{Modelos de comunicación (Punto 5):} JobFinder aplica un modelo cliente-servidor híbrido sobre TCP. Las operaciones transaccionales (registro, validación de ofertas RF-05 y postulación RF-04) usan REST síncrono; las alertas de vacantes (RF-06) y cambios de estado (RF-10) usan WebSockets asíncronos. Se descartó UDP porque la pérdida silenciosa de un paquete impediría notificar una entrevista al graduado. \textit{Práctica 2:} envío de mensaje simple (eco TCP) seguido del flujo bidireccional RF-10 (la empresa cita a entrevista y el graduado confirma asistencia en vivo).

\subsection{Bloque 3: Transparencias de Distribución y Teoremas CAP/PACELC (Minutos 6:00--9:00) -- Expositor: Heiner Ariel Zambrano Ronquillo}
\textbf{Transparencias (Punto 3):} JobFinder ofrece transparencia de acceso y ubicación al unificar los servicios detrás de rutas lógicas sin exponer puertos internos, y transparencia de concurrencia para múltiples sesiones WebSocket simultáneas. No garantiza transparencia total de fallos ante desconexiones abruptas de red móvil ni migración automática de sesiones en memoria.

\textbf{CAP y PACELC (Punto 4):} Ante una partición de red (\textbf{P}) hacia PostgreSQL, el sistema sacrifica disponibilidad de escritura y prioriza \textbf{Consistencia (C)} para evitar postulaciones duplicadas. Sin partición (\textbf{E}), PACELC equilibra baja \textbf{Latencia (L)} en notificaciones WebSocket y \textbf{Consistencia (C)} en transacciones relacionales. \textit{Práctica 3:} simulación de partición de base de datos con bloqueo controlado de escrituras.

\subsection{Bloque 4: Tolerancia a Fallos, Comparativa y Lecciones PFC (Minutos 9:00--12:00) -- Expositor: Luis Alberto Paca Pilataza}
\textbf{Tolerancia a fallos (Punto 6):} Se detectan enlaces caídos mediante tramas de control \texttt{PING/PONG} cada 5 segundos y \textit{timeouts} TCP, recuperando la sesión cliente mediante reconexión escalonada con \textit{backoff} exponencial (1\,s, 2\,s, 4\,s, 8\,s). Aplicando el análisis de causas raíz de FORM-AC-01 (caída de Facebook de 2021), se evita el acoplamiento entre procesos externos y el bucle principal de eventos.

\textbf{Comparación y lecciones PFC (Puntos 7 y 8):} Frente a LinkedIn y CompuTrabajo, JobFinder aporta validación administrativa previa (RF-05), asumiendo como compromiso (\textit{trade-off}) mayor consumo de RAM por mantener sockets TCP abiertos frente a APIs \textit{stateless}. Para la entrega \texttt{TA-PFC-E1} se aislará la importación externa (RF-07) en un contenedor Docker independiente con colas asíncronas. \textit{Práctica 4:} detención abrupta del servidor (\texttt{Ctrl+C}) y reconexión automática del cliente.

\end{document}